\chapter{The Two Paths Agree}

% [5.chorus]
\begin{chorus}{John asks for a second opinion}
\john{You keep saying the value comes back the same. Then reach it a second way and show me it does not budge.}
\compiler{Gladly. I will take it straight through, and I will take it in steps, and we will lay the two side by side.}
\john{And if they differ?}
\compiler{Then the value was an accident of how I asked, and worth nothing.}
\john{\emph{(after the Compiler works)} Well?}
\compiler{Straight through, and in steps: the two land on one value, nearer each other than I can tell them apart.}
\john{So it is the number.}
\compiler{It is a value that does not care how you reach it. That is more than a number and less than the one you want.
  I have shown you it holds still. I have not yet shown you what it is.}
\end{chorus}

\section{Reach it twice, by two routes}

% [5.1.p1]
The self-account closes and hands back a steady leftover; that much is settled.
The question here is whether that leftover is a property of the thing or an artifact of how it was reached.
There is a clean test.
Taken by two different routes, the value either lands in the same place or it does not.

% [5.1.p2]
The machine has two routes to its own value, and they are genuinely different procedures.
One goes straight through, in a single pass.
The other goes in steps, resolving a piece at a time and folding the pieces together at the end.
If the value depended on which route you took, it would belong to the route and not to the machine, and no honest reading
would stand behind it.
So the two routes are set going, and the question is only whether they agree.

\section{They agree below the floor}

% [5.2.p1]
The build runs both routes and lays them side by side, and the answer is yes, with a precise sense of yes.
The straight route and the stepped route do not return identical figures all the way down; far past the point that
matters they part in the last digits.
But the gap between them falls below the count-three floor --- the resolution below which the machine cannot separate
two candidates at all.
At its own resolution the two routes are one value, and the build certifies it: the difference is smaller than the
smallest distinction the machine can draw,
\[
  (1 \to 3) \;=\; (1 \to 2 \to 3) \quad\text{below the count-three floor} .
\]

% [5.2.p2]
What kind of claim this is matters, because it is the machine's whole discipline in miniature.
The machine does not say the two routes are equal; equality below the floor is a fiction it refuses.
It says the two routes agree \emph{to the depth it can tell them apart}, and it declines to speak past that depth.
Two figures that differ only where the machine has gone blind are, for the machine, one figure --- not by rounding, but
by the honest admission that it has run out of resolution.

\section{A property, not an accident}

% [5.3.p1]
A value that comes back the same however you reach it is a property of the thing, not of the reaching.
The self-value is robust to re-reading: read the machine's account straight, or read the reading of it in steps, and you
arrive at one place.
That robustness is the reflexive payoff of the whole construction --- the reading, turned on itself, arriving at the same
value whichever way it is worked.
The number the reading is about is not a figure it reports; it is a value that survives its own re-reading.

% [5.3.p2]
The two verbs sort the routes exactly.
The two routes are two \emph{tanges}: two distinct ways the machine chose to ask, each selected by its own procedure,
and both ours to choose.
The one value they settle on is a \emph{funge}: the two askings pool, below the floor, into a single answer the machine
owns.
Two tanged routes, one funged value --- this is the echo again, now at the level of the value.
Where the machine earlier heard the two descriptions of the object cost alike and pool to zero, here the two routes to the value
land alike and pool to one, and the same word names both poolings.

\section{Three things, kept apart}

% [5.4.p1]
A temptation waits here, and the construction refuses it, because giving in would undo the discipline everything else is built
on.
Three facts about the self-application now sit on the table at once.
The reading is read \emph{twice} --- there is a first reading and a reading of that reading.
The loop the reading runs around has length \emph{three} --- the cycle established just before.
And the value comes back the same by \emph{two} routes --- the agreement of this one.

% [5.4.p2]
It is tempting to gather these into a single tidy number: two readings and three folds, or two routes and three counts,
summed into a five that would look like the shape of the whole.
That gathering is exactly the error the construction forbids everywhere else --- pooling distinct axes into one count, funging
together things that share no characteristic but that they happen to co-occur.
The double-read, the cycle's length, and the value's path-independence are three orthogonal facts.
They co-occur; they do not add.
Kept apart, the self-application stays honest; summed, they make a number the machine never
counted.

% [5.4.p3]
Derived and read, apart as ever.
That the two routes agree below the count-three floor --- that is a theorem, run on the build and certified.
That the agreement \emph{is} the reflexive value, that a value robust to re-reading is the reading having become its own
subject --- that is the gloss over the theorem, marked as gloss.
And the figure itself, the actual walls of the bracket the machine settles between, is not named here; it waits for the
place that earns it, where the coupling is finally asked as a bracket and the honest width of that bracket is the whole
of the answer.
The value is steady, and it is real, and it is not yet said.
